\section{总结与延伸}

Lecture 12 的核心不是列举 benchmarks，而是建立评价思维：评价会把抽象目标压成具体指标，指标会改变模型开发方向。Perplexity 平滑但不真实；考试清晰但不开放；聊天偏好真实但有偏；agent benchmarks 有行动性但混入 scaffold；reasoning benchmarks 试图剥离知识但仍有先验；safety、realism、validity 则提醒我们，评价本身也是需要被评价的对象。

注意：本讲讨论的 zero-shot、zero contamination、one true evaluation 中的 zero 只是普通英文含义，不是 ZeRO（Zero Redundancy Optimizer）；ZeRO 是训练并行中的优化器状态、梯度和参数分片方法，Lecture 12 不以它为主题。

把本讲压缩成一个实践流程：先写下评价目的，再决定评价对象是 method、model 还是 agent；然后选择指标类型，是 probability、exam accuracy、human preference、tool-task success、safety refusal，还是 real-world workflow；接着检查 validity，包括污染、题目质量、judge bias、privacy 和 representativeness；最后才看 leaderboard 分数。这个顺序很重要，因为分数只有在规则清楚时才有含义。

因此，一份好的 evaluation report 不应该只列出模型名和分数，还应该说明数据来源、prompt 形式、是否允许工具、judge 是谁、失败样例长什么样、置信区间有多宽，以及这个评测不覆盖哪些真实场景。只有这些上下文齐全，读者才能判断一个分数是否值得信任。

\begin{importantbox}{最终 takeaways}
\begin{enumerate}
    \item 没有 one true evaluation；先说明评价目的和 rules of the game。
    \item Perplexity 仍有用，但不能代表真实 assistant quality。
    \item Exam benchmarks 易评分但会饱和，且不等同于真实使用。
    \item Open-ended chat 需要 preference、rubric 和 judge-bias 控制。
    \item Agent benchmarks 评价的是 model + scaffold + tools 的系统。
    \item Safety 和 realism 都高度上下文相关，隐私和公开复现之间有张力。
    \item Validity 包括 contamination、dataset quality、fresh/private evals 和 benchmark auditing。
\end{enumerate}
\end{importantbox}

\begin{knowledgebox}{拓展阅读}
\begin{itemize}
    \item HELM, MMLU, MMLU-Pro, GPQA, Humanity's Last Exam.
    \item Chatbot Arena, AlpacaEval, WildBench.
    \item SWE-Bench Verified, TerminalBench, CyBench, MLE-Bench.
    \item ARC-AGI, HarmBench, AIR-Bench, GDPVal, MedHELM.
    \item Train-test contamination and benchmark quality auditing papers.
\end{itemize}
\end{knowledgebox}

\end{document}
